\section{An RMS comparison for finite convex experiments}
Let $X$ be a finite-dimensional nonempty compact convex set and $\ell(x)$
a linear functional. Consider either of the following experiments:
\begin{enumerate}
\item independent finite categorical observations, whose positive probability
vectors $p_i(x)$ are affine functions of $x$;
\item a finite vector of independent Poisson counts, whose strictly positive
intensities $v_j(x)$ are affine functions of $x$.
\end{enumerate}
In the first case let $\mathcal F$ be the finite-dimensional space of
functions $c+\sum_i\phi_i(z_i)$. In the second let
$\mathcal F=\{c+\sum_j\beta_j z_j\}$. Let $P_x$ denote the whole
observation law. All exponential moments of functions in $\mathcal F$
are finite. For fixed $\phi\in\mathcal F$, the function
$x\mapsto\log\E_xe^\phi$ is concave: it is a sum of logarithms of positive
affine forms in the categorical case, and an affine function of the
intensities in the Poisson case. Moreover $\log(dP_x/dP_y)\in\mathcal F$.

\begin{lemma}[RMS affine comparison]\label{lem:comparison}
For any measurable estimator $G$ with
$r_G=\sup_{x\in X}(\E_x(G-\ell(x))^2)^{1/2}<\infty$ and any $\eta>0$,
there exists $F\in\mathcal F$ such that
\[
 \sup_{x\in X}(\E_x(F-\ell(x))^2)^{1/2}\le12r_G+\eta.       \tag{2}
\]
The result permits nonattainment of the infimum and imposes no restriction
on the original estimator's output range.
\end{lemma}
\begin{proof}
Put $s=\log2$. For $\alpha>0$ and $\phi\in\mathcal F$ define
\begin{align*}
 \Phi(x,y;\phi,\alpha)
 &=\ell(x)-\ell(y)
   +\alpha\log\E_y e^{\phi/\alpha}
   +\alpha\log\E_x e^{-\phi/\alpha}+2\alpha s,\\
 2S&=\inf_{\phi,\alpha>0}\max_{x,y\in X}\Phi(x,y;\phi,\alpha).
\end{align*}
This real-valued continuous function is convex in $(\phi,\alpha)$ by the
perspective transformation and concave in $(x,y)$ by the preceding
log-moment observation. Sion's minimax theorem applies because $X\times X$
is compact convex and the other domain is convex. Thus the infimum and
supremum commute. For fixed $x,y$, Cauchy--Schwarz gives
\[
 \inf_{\psi\in\mathcal F}
  \left(\log\E_y e^\psi+\log\E_x e^{-\psi}\right)
   =2\log\Aff(P_x,P_y),
\]
with equality at $\psi=\tfrac12\log(dP_x/dP_y)$.
Taking the infimum over $\alpha>0$ after putting $\phi=\alpha\psi$
therefore proves
\[
 2S=\max\{\ell(x)-\ell(y):x,y\in X,\ \Aff(P_x,P_y)\ge1/2\}. \tag{3}
\]
The constraint is closed and nonempty, so the maximum exists. Continuity
of affinity is immediate from the finite categorical formula and, for
Poisson vectors, from
$\Aff(P_x,P_y)=\exp[-\tfrac12\sum_j(\sqrt{v_j(x)}-\sqrt{v_j(y)})^2]$.
In particular $0\le S<\infty$.

For the maximizing pair in (3), total variation is at most $\sqrt3/2$
by the standard Cauchy--Schwarz affinity bound
$\|P-Q\|_{\rm TV}\le\sqrt{1-\Aff(P,Q)^2}$.
Testing the two endpoints by thresholding $G$ at their functional midpoint
gives, with $\Delta=2S$,
\[
 2r_G^2\ge \frac{\Delta^2}{4}
           (1-\|P_x-P_y\|_{\rm TV}),\qquad
 S\le\sqrt{\frac{2}{1-\sqrt3/2}}\,r_G.                    \tag{4}
\]
Indeed each wrong decision has squared estimation error at least
$\Delta^2/4$, and the sum of the two testing errors is at least one minus
total variation. The case $S=0$ requires no division by $S$.

Fix $R>S$. Choose $(\phi,\alpha)$ with
$\max_{x,y}\Phi\le2R$. On the diagonal, Cauchy--Schwarz gives
$\Phi(x,x;\phi,\alpha)\ge2\alpha s$, hence $\alpha\le R/s$.
Set
\[
 U=\max_x\{\ell(x)+\alpha\log\E_xe^{-\phi/\alpha}\},\qquad
 V=\max_x\{-\ell(x)+\alpha\log\E_xe^{\phi/\alpha}\}.
\]
Then $U+V+2\alpha s\le2R$. The additive estimator
$F=\phi+(U-V)/2$ satisfies, uniformly in $x$,
\[
 \E_x e^{(F-\ell(x))/\alpha}\le e^{R/\alpha-s},\qquad
 \E_x e^{(\ell(x)-F)/\alpha}\le e^{R/\alpha-s}.
\]
For $t\ge0$, Markov's inequality consequently gives
$\Pp_x(|F-\ell(x)|>R+t)\le e^{-t/\alpha}$, since $2e^{-s}=1$.
Integrating this tail, rather than substituting a confidence guarantee
for an RMS bound, yields
\[
 \E_x(F-\ell(x))^2
 \le R^2+2R\alpha+2\alpha^2
 \le R^2\left(1+\frac2s+\frac2{s^2}\right).                \tag{5}
\]
The product of the squared constants in (4)--(5) is less than $144$:
$\log2>2/3$ and $\sqrt3/2<7/8$ bound it by $8.5\cdot16=136$.
Letting $R\downarrow S$, with the slack chosen for the given $\eta$,
proves (2). This also proves uniform square integrability of the produced
estimator, rather than assuming it.
\end{proof}

\section{From a base-cost obstruction to one pair hard for every capped policy}
In this section only, restrict the parameter set to
\[
 \mathcal X_\rho=\{q:0\le q\le1/2,\ \int q\,d\mu\le\rho\},
 \qquad J=\rho L H.
\]
It is a weak-star compact convex subset of $L^\infty(\mu)$.
The following implication separates its one required entropy input from
the convex-experiment argument. The input is the base-cost bound proved in
the accompanying base-cost section of the full-model note:
\begin{equation}\label{eq:base-cost-input}
 \begin{gathered}
 \text{There is a universal }c_b>0\text{ such that every observable
 single-coin estimator }Y\text{ with}\\
 \sup_{q\in\mathcal X_\rho}|\E_qY-\theta(q)|\le\rho/64,
 \quad V=\sup_{q\in\mathcal X_\rho}\Var_qY<\infty,
 \quad \sup_{q\in\mathcal X_\rho}\E_qT<\infty
 \quad\text{satisfies}\quad V\E_0T\ge c_bJ.
 \end{gathered}                                                     \tag{BC}
\end{equation}
The cost in (BC) is at the all-negative population $q=0$, not the supremum
over populations. The two lemmas below are conditional on precisely this
input. No algorithm-level lower bound is inferred from a variational
optimization alone. Write
$h(P,Q)=1-\Aff(P,Q)$ for squared Hellinger distance in the normalization
used here; it takes values in $[0,1]$.

\subsection*{Exact finite observation reduction}
A capped local policy $A$ includes its independent seed and stops after at
most $m<\infty$ flips of one genuine coin. Its full observed transcript
contains the seed and the stopped word. Its seed distribution and policy
may depend on known parameters, but not on the actual $q$.
For each word $x\in\{0,1\}^t$, let $B_{t,x}$ be the measurable seed set
that causes stopping exactly at that word. For $0\le k\le t\le m$ put
\[
 a_{t,k}=\sum_{x:\,|x|=t,\,\#1(x)=k}\Pp(\text{seed}\in B_{t,x}),
\]
and, with $p_s(d)=1/2+(2s-1)d$, put
\[
 m_{t,k}(q)=\int\big[(1-q(d))p_0(d)^k(1-p_0(d))^{t-k}
                  +q(d)p_1(d)^k(1-p_1(d))^{t-k}\big]\,d\mu.
\]
The finite statistic $(T,K)$, where $K$ is the number of observed heads,
has probabilities $a_{t,k}m_{t,k}(q)$. Conditional on $(T,K)$, the stopped
word and seed have a distribution independent of $q$: their common
likelihood multiplier is exactly $m_{t,k}(q)$ and cancels. This statement
also follows directly by integrating any bounded transcript function over
the finitely many word/seed cells, so it does not assume a latent oracle.
Consequently the full transcript and this finite statistic have exactly
the same Hellinger affinities for all pairs of populations.

Remove coordinates with $a_{t,k}=0$. The remaining coordinates are positive
for every $q$, since $p_s(d)\in[1/4,3/4]$. They are affine weak-star
continuous functions of $q$. Thus their probability vector together with
$\theta(q)$ has a compact convex finite-dimensional range. In particular,
for each capped $A$, its full-transcript squared Hellinger distance
$h_A(q_-,q_+)$ is a continuous convex function of the pair on
$\mathcal X_\rho^2$. Convexity follows from joint concavity of affinity for
two affine probability vectors. This reduction retains the stopping time
and expected cost, not just a fixed number of sampled bits.

\begin{lemma}[A hard pair for each capped policy]\label{lem:hard-pair}
Assume (BC). There are universal constants $c_\delta>0$ and $C_H<\infty$
such that, for every $\rho\in(0,1/4]$, every
$0<\delta\le c_\delta\rho$, and every randomized capped local policy $A$,
there exist $q_-,q_+\in\mathcal X_\rho$ with
$\theta(q_+)-\theta(q_-)\ge\delta$ and
\[
 h_A(q_-,q_+)\le C_H\frac{\delta^2}{J}\E_0T_A.              \tag{6}
\]
The full observed seed and stopped history are included in the Hellinger
distance.
\end{lemma}
\begin{proof}
Let $w=\E_0T_A$. If $w=0$, the seed-only event of taking a first observation
has probability zero, so $A$ has the same transcript law for every $q$.
The pair $q_-=0$, $q_+=\delta$ gives separation $\delta$ and $h_A=0$.
We therefore assume $w>0$ in all the divisions below.
First establish the following additive lower bound for
$n\ge1$ independent runs of $A$. If an additive output with bounded
component functions
$F=c+\sum_{i=1}^nY_i$ has worst RMS $r<\rho/64$, put
$b_i=\E_0Y_i$, $b=\E_0F$, and choose $I$ uniformly in $\{1,\ldots,n\}$
independently. Run just its component on one coin and return
$Z=n(Y_I-b_I)+b$. This preserves every population mean and uses base cost
$w$.

For $\theta=\theta(q)$, the laws of each component output under $q$ and
$0$ have a common subprobability of mass $1-\theta$: integrate the
negative-sign kernel with weight $1-q(d)$. For any two distributions
sharing such a subprobability, writing each as a mixture of its normalized
common part and residual part shows
\[
 (\E_qY_i-\E_0Y_i)^2
 \le\frac{2\theta}{1-\theta}
       (\Var_qY_i+\Var_0Y_i).                              \tag{7}
\]
Indeed the squared displacement of either mixture mean from the common
mean is at most $\theta/(1-\theta)$ times its variance; use
$(u-v)^2\le2u^2+2v^2$. The case $\theta=0$ is identical laws.
Independence and $\sup_q\Var_qF\le r^2$ imply
\[
 \sum_i(\E_qY_i-b_i)^2\le\frac{4\rho}{1-\rho}r^2,
 \quad
 \Var_qZ\le n\sum_i\Var_qY_i
                 +n\sum_i(\E_qY_i-b_i)^2
              \le\frac73nr^2.
\]
All outputs here are bounded, and stopping is capped, so (BC) applies to
$Z$ without an additional integrability assumption. It gives
$r^2\ge3c_bJ/(7nw)$. Including the branch $r\ge\rho/64$ establishes
\[
 r\ge c_a\min\{\rho,\sqrt{J/(nw)}\},\qquad
 c_a=\min\{1/64,\sqrt{3c_b/7}\}.                           \tag{8}
\]

For $w>0$, let $S_n$ be the saddle quantity $S$ from the proof of
Lemma~\ref{lem:comparison} for $n$ independent copies of the finite
statistic of $A$, on the parameter set $\mathcal X_\rho$.
Equations (3) and (5) of that proof show
\[
 2S_n=\max\{\theta(q_+)-\theta(q_-):
             \Aff(P_{A,q_+},P_{A,q_-})^n\ge1/2\},
 \quad \inf_{F\text{ additive}}\sup_q\|F-\theta(q)\|_2\le3S_n.
\]
The latter constant is valid because
$1+2/\log2+2/(\log2)^2<9$. The finite statistic outputs are bounded, so
(8) applies. Hence, for $\beta=2c_a/3$,
\[
 2S_n\ge\beta\min\{\rho,\sqrt{J/(nw)}\}.                 \tag{9}
\]
Take $c_\delta=\beta/2$ and $C_H=8/\beta^2$.
For $0<\delta\le c_\delta\rho$ set
$x=\beta^2J/(4\delta^2w)$. If $x\ge2$, choose $n=\lfloor x\rfloor$,
so $x/2\le n\le x$. Equation (9) gives a pair of separation at least
$2\delta$, with
\[
 h_A\le1-e^{-(\log2)/n}
       \le\frac{\log2}{n}
       \le\frac{8\log2}{\beta^2}\frac{\delta^2w}{J}
       \le C_H\frac{\delta^2w}{J}.
\]
If $x<2$, the pair $q_-=0$, $q_+=\delta$ belongs to
$\mathcal X_\rho$, has separation $\delta$, and $h_A\le1$ suffices
because $C_H\delta^2w/J>1$.
\end{proof}

\begin{lemma}[One pair hard for every capped policy]\label{lem:universal-pair}
Assume (BC), and keep the universal constants of
Lemma~\ref{lem:hard-pair}. For every $0<\rho\le1/4$ and
$0<\delta\le c_\delta\rho$, there exist
$q_-,q_+\in\mathcal X_\rho$, with
$\theta(q_+)-\theta(q_-)\ge\delta$, such that simultaneously for every
randomized capped single-coin policy $A$,
\[
 h_A(q_-,q_+)\le C_H\frac{\delta^2}{\rho LH}\E_0T_A.       \tag{10}
\]
Each policy's seed and full stopped observable history are included.
The constants and pair are chosen before the policy.
\end{lemma}
\begin{proof}
Let $\mathcal K_\delta=\{(q_-,q_+)\in\mathcal X_\rho^2:
\theta(q_+)-\theta(q_-)\ge\delta\}$. It is nonempty compact convex in
the weak-star topology. Fix any finite list of capped policies
$A_1,\ldots,A_k$ and put $w_i=\E_0T_{A_i}$ and
$\lambda=C_H\delta^2/(\rho LH)$. For a simplex vector
$\pi\in\Delta_k$, define the observable mixture policy that draws
$i\sim\pi$, reveals this index in its independent seed, and runs $A_i$.
This is still capped, by the largest of the finitely many caps. Its
base cost is $\sum_i\pi_iw_i$ and, since the revealed labels occupy
disjoint transcript components, its squared Hellinger distance is exactly
$\sum_i\pi_i h_{A_i}$. Lemma~\ref{lem:hard-pair} therefore gives
\[
 \min_{z\in\mathcal K_\delta}
       \sum_i\pi_i\big(h_{A_i}(z)-\lambda w_i\big)\le0
       \qquad(\pi\in\Delta_k).
\]
For the finite simplex the displayed function is continuous affine in
$\pi$ and continuous convex in $z$, by the finite observation reduction.
Sion's theorem now gives
\[
 \min_{z\in\mathcal K_\delta}
       \max_{1\le i\le k}\big(h_{A_i}(z)-\lambda w_i\big)
   =\max_{\pi\in\Delta_k}\min_{z\in\mathcal K_\delta}
       \sum_i\pi_i\big(h_{A_i}(z)-\lambda w_i\big)\le0.
\]
Thus the closed subsets
$\{z\in\mathcal K_\delta:h_A(z)\le\lambda\E_0T_A\}$
have the finite intersection property as $A$ ranges over all capped
policies. Compactness of $\mathcal K_\delta$ gives a point in their full
intersection, proving (10). This uses only finite revealed mixtures and
a finite simplex minimax theorem; no topology on a space of all policies,
unjustified exchange of arbitrary infima, or unobserved-mixture affinity
identity is needed.
\end{proof}

